\scanpage{93}
% archive foliation: 92
On sait que sa composée avec
\[
\mathcal{L}\otimes\underline{\omega}\longrightarrow\underline{\omega}
\]
est $\unsure{-2i_{\underline{\omega}}|-2\operatorname{id}_{\underline{\omega}}}$. \deleted{\unsure{Considérons}} \unsure{l'hom.} $c|\Omega$.
$c|\underline{\Omega}$ est op. diff d'ordre $\leq1$, on sait qu'il se factorise en
\diagram{\begin{tikzcd}
\underline{\Omega}\arrow[r,"{d^1_{\unsure{\omega}}}","{\substack{\text{op. diff.}\\\text{can.}\\\text{d'ordre }1}}"'] & P^1_{\unsure{X/S}}(\underline{\Omega})\arrow[r,"\gamma","{\uline{\text{lin.}}}"'] & \mathcal{L}\otimes\underline{\omega}
\end{tikzcd}}
et comme la composée
\diagram{\begin{tikzcd}
\underline{\Omega}\arrow[r,"{c|\underline{\omega}}"] & \mathcal{L}\otimes\underline{\omega}\arrow[r] & \underline{\omega}
\end{tikzcd}}
est linéaire (de valeur $-2i$), on trouve que $\gamma$ donne lieu à une hom. d'extensions
\diagram{\begin{tikzcd}
0\arrow[r] & \Omega\otimes\omega\arrow[r]\arrow[d,"{\gamma_0=?}"] & P^1_{X/S}(\Omega)\arrow[r]\arrow[d,"\gamma"] & \Omega\arrow[r]\arrow[d,"{-2i}"] & 0\\
0\arrow[r] & \Omega\otimes\omega\arrow[r] & \mathcal{L}\otimes\omega\arrow[r] & \omega\arrow[r] & 0
\end{tikzcd}}

On va déterminer $\gamma_0$, \unsure{quand on a} splitting $s$ et section $\delta_0$ de $\Omega$, en termes de $\unsure{\omega_0=i|\omega_0=\delta}$, $\omega_1$, $\omega_2$ ci-dessus, \deleted{\unsure{on voit que}} \inserted{En vertu de la} \unsure{recette générale valable pour les op.} $\leq1$ diff.\deleted{\unsure{$P^1_{X/S}(\Omega)\simeq P^1_{X/S}(\mathcal{O}_X)$ et les}}\inserted{sur un $\mathcal{O}_X$-mod. $\Omega$,}
\margin{$\unsure{\delta_0\text{ de }\Omega}$}
ainsi
\[
\gamma(\underset{\in\Omega}{\omega_1}\otimes d\underset{\in\mathcal{O}_X}{f})=c(f\omega_1)-f\,c(\omega_1).
\]
Si $\omega_1=\lambda\delta$, on trouve donc
\[
\begin{aligned}
\gamma(\lambda\delta\otimes df)&=\delta\,[d(f\lambda)-f\,d(\lambda)]\\
&=\delta\,[\lambda\,df]=\lambda\delta\otimes df.
\end{aligned}
\]
donc $\gamma=\operatorname{Id}$. Donc

\uline{Prop.} Si $c$ connexion \uline{transverse}, \inserted{et $2$ inv. sur $X$,} alors $\gamma$ est un \uline{iso} d'extensions de $P^1_{X/S}(\Omega)$
